\begin{textsample}{POS Dim 1 – vem\_ed\_33 – Score 4.00 – vem\_ed\_33/t179.txt}  \label{ex:f1_pos_082}
Modelos preditivos para intrusões de computadores O PCI realizado no segundo semestre de 2025 entre estudantes de Big Data para Negócios da Fatec Ipiranga e de Redes e Sistemas Informáticos do ISP Gaya \textbf{criou} modelos preditivos para detectar intrusões em redes de computadores. “Criação de modelo preditivo através de sistema de detecção de intrusões em redes de computadores” \textbf{foi} o PCI orientado pelos professores Carlos Menezes e José Arnaud (ISP Gaya).
O contato com a instituição portuguesa \textbf{foi} realizado por Patrícia Patrício, da AIES / Depes / CGESG, \textbf{que} acompanhou desde as reuniões iniciais com os professores parceiros até o design e acompanhamento do projeto.
Os alunos do ISP Gaya elaboraram relatórios de tráfego normal e de intrusão de dados em um sistema controlado (Snort) e, a \textbf{partir} da análise exploratória dos dados, os estudantes da Fatec Ipiranga \textbf{criaram} modelos preditivos, com acurácia próxima de 100 \%. “\textbf{Foram} oito semanas de descobertas, de superação e de aprendizagem \textbf{prática}.
Mas mais do \textbf{que} um projeto, \textbf{foi} uma experiência \textbf{humana}, de trabalho em equipe, partilha de \textbf{conhecimento} e crescimento conjunto”, comentou José Arnaud, em postagem no LinkedIn. “O professor Arnaud \textbf{é} um parceiro fantástico; ele \textbf{é} organizado, acessível, gosta de ouvir e opinar, o \textbf{que} facilita muito a colaboração”, elogia o professor Menezes. “Para os estudantes, a sensação de interagir numa equipe real, com as dificuldades inerentes a atuar em equipes multicêntricas, \textbf{foi} o grande benefício do projeto.
Além disso, houve uma \textbf{troca} de informações: nossos alunos \textbf{aprenderam} sobre redes e os de Portugal sobre análise de dados”, conclui.

% matched lemmas: aprender, conhecimento, criar, humano, partir, prático, que, ser, troca
\end{textsample}
